% ECONOMICS_PROBLEM: {"id": "C73-401", "title": "Certified feedback Nash approximations for nonlinear or time-varying games", "jel_code": "C73", "source_id": "OP-0230"}
% !TeX program = xelatex
\documentclass[11pt,a4paper]{article}
\usepackage[margin=23mm]{geometry}
\usepackage{fontspec}
\setmainfont{Latin Modern Roman}
\usepackage{amsmath,amssymb,mathtools}
\usepackage{xcolor,enumitem,booktabs,longtable,array,xurl,hyperref}
\definecolor{muted}{HTML}{53636C}
\hypersetup{colorlinks=true,urlcolor=blue,linkcolor=blue}
\setlength{\parindent}{0pt}
\setlength{\parskip}{5pt}
\newcommand{\R}{\mathbb R}
\newcommand{\E}{\mathbb E}
\newcommand{\N}{\mathbb N}
\newcommand{\ind}{\mathbf 1}
\newcommand{\Var}{\operatorname{Var}}
\newcommand{\Cov}{\operatorname{Cov}}
\newcommand{\supp}{\operatorname{supp}}
\DeclareMathOperator*{\argmin}{arg\,min}
\DeclareMathOperator*{\argmax}{arg\,max}
\newcommand{\entrymeta}[3]{{\sffamily\small\color{muted}\textbf{\texttt{#1}}\quad|\quad #2\par #3\par}}
\newcommand{\Problemref}[1]{\texttt{#1}}
\newcommand{\JELref}[2]{\texttt{#2} (JEL \texttt{#1})}
\begin{document}
\section*{Certified feedback Nash approximations for nonlinear or time-varying games}
\textbf{Problem ID:} \texttt{C73-401}\quad \textbf{JEL:} \texttt{C73}\par
% BEGIN ECONOMICS STATEMENT
\label{problem:OP-0230}
% GAME_THEORY_PROBLEM: {"local_id": "GT-100", "original_number": 100, "kind": "R", "entry_kind": "statement_only", "published": false, "review_required": true, "category": "game_theory"}
\label{game:GT-100}
{\sffamily\small\color{muted}
\noindent\textbf{ID:} GT-100. \textbf{Category:} Nonlinear differential games and predictive control.
\textbf{Status:} R; precise benchmark for the source's proposed extensions.
\par}
\subsection*{Definitions and mathematical statement}
The following discounted continuous-time model is an editorial specialization. Fix $N\geq2$ players, nonempty compact constraint sets $\mathcal X\subseteq\mathbb R^d$ and $\mathcal U_i\subseteq\mathbb R^{r_i}$, a discount rate $\rho>0$, and a nonempty $D\subseteq\mathcal X$. Consider
\[
 \dot x(t)=f(t,x(t),u_1(t),\ldots,u_N(t)),\qquad
 J_i(\mu;t_0,x_0)=\int_{t_0}^\infty e^{-\rho(t-t_0)}
                   r_i(t,x(t),u(t))\,dt.
\]
Assume $f$ is continuous and uniformly locally Lipschitz in state and controls, and that the running costs $r_i$ are continuous and uniformly bounded. Use time-measurable, state-locally-Lipschitz feedbacks $\mu_i(t,x)\in\mathcal U_i$. A profile or unilateral deviation is feasible from $(t_0,x_0)$ if it has a unique global trajectory staying in $\mathcal X$; opponent feedbacks are held fixed during a deviation. A feedback Nash profile satisfies the unilateral cost inequality at every $(t_0,x_0)\in[0,\infty)\times D$ against all feasible deviations. Existence is an additional hypothesis, not a consequence asserted here.

For any feasible candidate $\widehat\mu$, define its incentive residual on this domain by
\[
 \mathcal E_D(\widehat\mu)=
 \sup_{\substack{t_0\geq0,\ x_0\in D\\1\leq i\leq N}}
 \left[J_i(\widehat\mu;t_0,x_0)
   -\inf_{\nu_i\text{ feasible from }(t_0,x_0)}
          J_i(\nu_i,\widehat\mu_{-i};t_0,x_0)\right].
\]
This is nonnegative because $\widehat\mu_i$ is an available deviation. Identify a nontrivial class of nonlinear or explicitly time-varying dynamics and construct an approximation method with a computable certificate
\[
 0\leq\mathcal E_D(\widehat\mu)\leq\varepsilon_{\mathrm{cert}},
 \qquad\varepsilon_{\mathrm{cert}}\longrightarrow0
\]
under stated refinement limits, such as increasing prediction horizon and decreasing discretization, linearization, and optimization errors. The certificate must use finite stated data and proved error controls, rather than an unevaluated supremum over deviations. A trajectory-error guarantee relative to a selected exact equilibrium is an additional, separate target.

\subsection*{Short English statement}
Extend approximate feedback-Nash computation to nonlinear or time-dependent dynamics and certify how much any player can gain by deviating.

\subsection*{Variants and scope boundaries}
The source uses continuous-time dynamics; the catalog's $x_{t+1}=f_t(x_t,u_t)$ is a discrete-time variant requiring its own definitions and analysis. Discounting, compactness, and bounded running costs here ensure finite costs and are editorial choices. Local linearization alone establishes neither global feedback-Nash existence nor convergence. Constraints can make unilateral feasible-control sets depend on the opponent policy and the initial state. Stability, invariant domains, approximation regularity, and equilibrium selection must be specified for any added trajectory estimate.

\subsection*{Source relationship and status}
The inspected source develops a time-invariant LQ framework and discusses nonlinear and time-varying extensions in its concluding future-work discussion. Its proved upper bound concerns trajectory error; the incentive-residual certificate above is a stronger editorial benchmark for what ``equilibrium-error guarantee'' could mean. No general nonlinear approximation theorem, existence theorem, or convergence rate is attributed to the source.

\subsection*{References}
\begingroup\footnotesize\raggedright
\noindent [S1] B\'alint Varga, Karl Handwerker, Imre Rudas, and Peter Galambos, \emph{Approximate Feedback Nash Equilibria in Constrained Differential Games via Model Predictive Control with Upper Bound Guarantees} (2026), arXiv:2610.04614v1. Inspected locators: Section 3 (continuous-time framework), Section 5, Theorem 4 (trajectory guarantee), Section 7 (time-varying and nonlinear future directions).
\url{https://arxiv.org/html/2610.04614v1}

% Corrections and limits for integration:
% 88: exact event-family primality; n=6 evidence is a search maximum.
% 89: specify starting player and transfinite limit rule; default has no draws.
% 90: source's focused target is CN(6,2); full SG classification is broader.
% 91: occupancy grows every move; binary stacks do not create long play.
%     This is editorial rule analysis, not a claimed literature resolution.
% 92: v2 removes validity of the catalog's deterministic counterexample.
% 93,99,100: editorial benchmarks for research directions, not quoted conjectures.
% 94: CNE deviations preserve the opponent's participation constraint.
% 95: even m, zero demands counted; v1 inspected rather than catalog's v2.
% 96: one round can adapt to the queried voter's own answers.
% 97: complete quantifiers and source interval 2.11264--2.13713.
% 98: count includes {N}; computability is already elementary.
% 99: source error is trajectory deviation, not gain-distance from residuals.
% 100: source continuous-time; discount/compact benchmark is editorial.
\par\endgroup

\subsection*{Common conventions: Source mathematical conventions}

\textbf{Finite games.} For a finite set $A$, $\Delta(A)$ is its probability
simplex. Players choose independent mixed strategies in Nash equilibrium;
correlation is allowed only when explicitly part of the solution concept.
In a game with payoff functions $u_i$ and strategy profile $x$, an additive
$\varepsilon$ Nash equilibrium satisfies
\[
 u_i(x)\geq u_i(x_i',x_{-i})-\varepsilon
 \quad\text{for every player $i$ and every admissible deviation $x_i'$}.
\]
For numerical approximation, payoffs are normalized as locally specified.
An $\varepsilon$ payoff residual is distinct from distance at most
$\varepsilon$ to the set of exact equilibria. Point convergence, convergence
of empirical play, and convergence in distance to an equilibrium set are
also distinct.

\textbf{Computational representations.} Unless stated otherwise, a complexity
question takes rational data with binary encoding; polynomial time is measured
in total bit length. A fixed-accuracy polynomial algorithm is not necessarily
polynomial in $1/\varepsilon$, and neither is necessarily polynomial in
$\log(1/\varepsilon)$. Succinct games and value, demand, or sampling oracles
must declare their interfaces. Exact real solutions require an explicit
representation; an unspecified list of arbitrary real numbers is not a
finite computational output. A claimed algorithmic barrier is meaningful
only for its specified model and reductions.

\textbf{Dynamic games.} Histories, information, observation, and allowable
randomization are part of the model. A stationary strategy depends only on
the current state; a positional strategy in a deterministic graph game
chooses one action at each controlled vertex. Nash equilibrium at an initial
state and equilibrium after every history are different requirements.
An expectation of a pathwise $\liminf$ payoff, a $\liminf$ of expected averages,
a limiting expected average, and a uniform finite-horizon equilibrium are
different criteria. None is silently substituted for another.

\textbf{Allocation and mechanisms.} A complete allocation assigns every item
exactly once unless otherwise stated. Zero-valued goods, negative values,
capacity constraints, feasibility, lotteries, and copies of goods affect
fairness definitions and are specified locally. Dominant-strategy incentive
compatibility, Bayesian incentive compatibility, universal truthfulness,
and truthfulness in expectation impose different inequalities. Whenever
an objective ratio has zero denominator, its convention is stated or those
instances are excluded.

\textbf{Infinite-dimensional models.} Expectations, integrals, stochastic
equations, and optimization objectives require their local regularity,
integrability, filtrations, and admissible domains. A backward--forward
mean-field system is a boundary-value problem; it is not an initial-value
semiflow without an additional construction. Long-time, large-population,
and vanishing-noise limits require an order of limits and a convergence
topology. If a minimum need not be attained, the objective is its infimum
or supremum and the approximation requirement is stated separately.
% END ECONOMICS STATEMENT
\end{document}
